\chapter{Introduzione}

La sicurezza informatica nei sistemi distribuiti è un campo in continua evoluzione, spinto dalla crescente complessità delle infrastrutture moderne e dall'inadeguatezza dei modelli di protezione tradizionali. Il paradigma del perimetro difensivo, in cui tutto ciò che si trova all'interno della rete aziendale è considerato affidabile per definizione, ha mostrato i propri limiti in modo sempre più evidente negli ultimi anni. Violazioni sofisticate, movimenti laterali non rilevati e compromissioni di credenziali legittime hanno reso chiaro che la fiducia implicita nella rete interna è una premessa insostenibile.

Da questa consapevolezza nasce il modello Zero Trust, la cui idea fondante può essere riassunta in un unico principio: non fidarsi mai, verificare sempre. In un'architettura Zero Trust, nessuna richiesta di accesso viene considerata sicura a priori, indipendentemente dalla provenienza, dall'identità dichiarata o dalla rete di appartenenza. Ogni accesso deve essere autorizzato esplicitamente, valutando in modo continuo il contesto dell'identità richiedente, del dispositivo utilizzato, della risorsa a cui si vuole accedere e del livello di rischio associato alla transazione.

\section{Il progetto}

Questo elaborato descrive la progettazione e l'implementazione di un'architettura Zero Trust basata su microservizi, sviluppata nell'ambito del corso di \textit{Advanced Cybersecurity for IT}. Lo scenario scelto come dominio applicativo è quello di una centrale nucleare, un contesto volutamente estremo per quanto riguarda i requisiti di sicurezza: ruoli operativi con responsabilità nettamente distinte, informazioni classificate a vari livelli di sensibilità, accesso fisico e digitale strettamente correlati e impatti potenzialmente critici in caso di compromissione.

L'architettura è distribuita su 7 componenti principali:
\begin{enumerate}
    \item PEP (Envoy PROXY)
    \item PDP (Security Orchestrator, OPA, Modello AI, Security DB)
    \item IAM 
    \item Logica (Business Logic, Business DB)
    \item Firewall
    \item IDS 
    \item SIEM (Splunk)
\end{enumerate} 
conforme ai principi definiti dalla pubblicazione NIST Special Publication 800-207, il riferimento tecnico consolidato per le implementazioni Zero Trust. Il fulcro del progetto è la separazione netta tra il Policy Enforcement Point (PEP), responsabile dell'intercettazione del traffico e dell'applicazione delle decisioni, e il Policy Decision Point (PDP), che valuta le policy e produce il verdetto di accesso. Questa separazione garantisce che nessun componente concentri in sé sia la capacità di applicare le regole sia quella di definirle, riducendo la superficie di attacco e aumentando la chiarezza architetturale.

Il sistema implementa un controllo degli accessi adattivo e basato sul rischio, che va ben oltre la semplice validazione delle credenziali. Ogni richiesta viene analizzata tenendo conto dell'impronta TLS del client tramite JA3 fingerprinting, della presenza di un certificato client valido mediante mutual TLS, dello score di anomalia restituito dal modello di ai e dalle policy di accesso definite. Il risultato è un sistema capace di prendere decisioni di accesso granulari, contestuali e dinamiche.

\section{Obiettivi}

Gli obiettivi del progetto sono articolati su più livelli. Sul piano architetturale, l'obiettivo è dimostrare come i principi Zero Trust possano essere applicati concretamente in un'infrastruttura containerizzata, con una segmentazione di rete rigorosa e una catena di autorizzazione verificabile end-to-end. Sul piano della sicurezza, si vuole mostrare come sia possibile implementare un controllo degli accessi che tenga conto simultaneamente di identità, dispositivo, rete, comportamento e classificazione della risorsa, superando i limiti dei modelli basati esclusivamente su ruoli o credenziali statiche.

Sul piano dell'osservabilità, un obiettivo esplicito è garantire la tracciabilità completa di ogni richiesta attraverso tutti i componenti del sistema, mediante la propagazione di un identificatore di correlazione univoco (\texttt{X-Request-ID}) e la centralizzazione dei log su Splunk. Tale integrazione consente sia di ricostruire il percorso di ogni transazione, sia di rilevare scenari di attacco complessi attraverso interrogazioni e allarmi configurati sulla piattaforma di monitoraggio.


\section{Struttura dell'elaborato}

L'elaborato è organizzato in modo da accompagnare il lettore dalla teoria alla pratica in modo progressivo. Il Capitolo 2 fornisce il quadro teorico completo, trattando i fondamenti della crittografia applicata, i modelli formali di controllo degli accessi, i principi di progettazione sicura e le tecnologie specifiche impiegate nel progetto. Il Capitolo 3 descrive l'architettura del sistema nei suoi componenti, flussi e scelte progettuali. Il Capitolo 4 entra nel dettaglio implementativo di ciascun modulo. Il Capitolo 5 presenta il modello di Intelligenza Artificiale di tipo Temporal Graph Network impiegato per la valutazione del rischio, descrivendone l'architettura, l'addestramento e i risultati ottenuti. Il Capitolo 6 presenta i test di carico e di sicurezza eseguiti sul sistema. Infine, il Capitolo 7 riporta le conclusioni dell'elaborato.
